% Transcription — fonds Alexandre Grothendieck, shelfmark 57, batch 11 (pages 201-220).
% Established from the facsimile archives/batches/57/batch-11.pdf, per the skill
% transcribe-grothendieck.
%
% Pass: Opus 5.5 (claude-opus-5-5), 2026-10-03 — first pass, unchecked against the pages by a human.
%
% Pages transcribed: 201-205, 207-211, 213-220. Absent: 206 (blank but for
% show-through), 212 (blank).
%   201-205: neutral component G_0 and torsion subgroup G^tau of a group scheme
%     locally of finite type over a field; Gamma = G/G_0, nGamma, Gamma^tau.
%   207-208: criterion for a morphism (or constructible part) to be quasi-compact
%     via proper fibres and the generic point.
%   209-211: constructibility/quasi-compactness of the prime-to-p part P' of Pic^tau.
%   213-218: example xy = t z^2 over k[t]: Pic of the degenerate conic, non-
%     representability by a scheme of finite type, the non-separated S_I.
%   219-220: Pic of the cone over X; family of quadrics degenerating to a cone.
\documentclass[11pt,a4paper]{article}
\input{../preamble/grothendieck}

\folder{57}
\batch{11}
\pages{201}{220}
\foldertitle{Picard : tapuscrit annoté (s.d.), lettres (s.d., 1962).}
\dating{1962-[vers 1968]}
\watermark{Édition de démonstration}

\begin{document}

\section{Composante neutre et sous-groupe de torsion d'un schéma en groupes}

\page{201}
\note{la page commence au milieu d'une démonstration commencée avant ce lot : il s'agit de montrer que la composante neutre $G_0$ d'un schéma en groupes $G$ localement de type fini sur un corps $k$ est de type fini et irréductible}
\ill{} quasi-compact de type fini sur $k$. Soit en effet $x$ un pt de $G_0$ rat.\ sur $k$, il suffit de prouver que $x \in U^2$ [N.B. $U^2$ est constructible \uncertain{quasi-compact} dans $G_0$] i.e.\ \struck{$xU \cap U$} l'existence de $y$ et $z \in U(k)$ tels que $yz = x$, i.e.\ $y = xz^{-1}$ ; ce qui signifie que $(U \cap xU^{-1})(k) \neq \emptyset$, \uncertain{or} $U$ et $xU$ sont des ouverts non vides dans $G_0$, \struck{\ill{}}. Donc $G_0$ est de type fini sur $k$. S'il n'était irréductible, tt point fermé de $G_0$ \uncertain{appartiendrait} à \struck{\ill{}} deux composantes irréductibles distinctes (par translation). Absurde. Le cas \ill{} \struck{\ill{}} 3° \uncertain{pour} \ill{} de \struck{considérons les propriétés $G_k$} structure [il suffit de considérer les composantes connexes des pts fermés].

\struck{\ill{}} Prenant $U = G_0$, \ill{} d'où \ill{} 1°. \struck{\ill{}} $G_0$ est un sous-schéma \uncertain{en groupes}. Il est invariant \add{\ill{} \uncertain{groupe}} \struck{Prenant} \uncertain{Soit} le \uncertain{quasi-longueur}. \struck{Soit} Il suffit de déterminer les composantes connexes $U$ des pts fermés $x$. \struck{\ill{}}

\struck{\ill{} \ill{} \ill{} \ill{} \ill{}}
[\ill{} $k'$, donc les composantes connexes \ill{} \ill{} \ill{}]

\drawing{une courbe ondulée, à trois bosses, tracée en travers des dernières lignes encadrées ; elle figure sans doute les composantes connexes de $G_{\bar k}$}

\page{202}
$\mathcal{H}$ \ill{} \ill{} \ill{} \ill{} fini et plat \uncertain{de}

$G_{\bar k}$ \ill{} dans \ill{}, \struck{soit $U'$} la \uncertain{réunion} des
\ill{} leurs composantes \ill{} \ill{} $U_{\bar k}$ [\uncertain{réf}]
\struck{donc $U$ \ill{}} $\bar k$, donc $U$ \uncertain{défini} \ill{} de type fini, sur $k$. Comme $U_{\bar k}$ est
loc.\ irréductible, il \ill{} \ill{} de \ill{} \ill{} $U$ [\uncertain{réf}]. Comme $e$ est rat./$k$, il s'ensuit \ill{} \ill{} \add{\ill{}} \ill{} défini \ill{} \ill{}
[\uncertain{cf}] \ill{} \ill{} ! irréductibilité ; \struck{\ill{}} \add{Comme} $G_{\bar k}$ est irréductible, $G_0$ est irréductible.

Pour \ill{} l'existence du \uncertain{quotient}, \ill{} \ill{} \ill{} \ill{}

\textbf{Corollaire} \quad $G$ \uncertain{schéma en groupes} loc.\ de type fini sur $k$, \ill{} \ill{} \uncertain{schéma en groupes} \ill{}, \struck{$H$ un sous-groupe invariant, donc} $G/H$ existe, et $G \to G/H$ est fid.\ plat (\uncertain{cas} où $H$ est distingué) \struck{\ill{} \ill{} \ill{} \ill{} $H = G_0$ :} Dans le cas $H = G_0$, \ill{} composante \ill{} de $G$, \ill{} $G_0$ [\ill{} \uncertain{signaler} …] et $G/G_0$ existe \struck{\uncertain{groupe}}, \ill{} \ill{} \ill{} \uncertain{étale} \ill{} fini sur $k$.

\marginal{\ill{} \uncertain{particulier} : \ill{} \ill{} \ill{}, \ill{} \ill{} \ill{} $H$ \ill{} \uncertain{fid.\ plat} \ill{} $G$ \ill{} \ill{}}

\page{203}
Le cas $H$ \uncertain{quelconque} s'en déduit.

\struck{Note} Autrement \uncertain{dit} : $G$ \uncertain{simple} sur $k$ $\Longleftrightarrow$ $G$ séparable sur $k$ $\Longleftrightarrow$ $G$ séparable sur $\bar k$ \ill{}

Supposons $G$ \uncertain{commutatif}, \struck{soit} donc \add{(séparable sur $k$, commutatif)} $\Gamma = G/G_0$ est un \struck{\ill{}} \add{\emph{discret}} $k$-groupe commutatif, \ill{} \ill{} \ill{} \ill{}. Soit $\Gamma \xrightarrow{\ n\ } \Gamma$ \ill{} \ill{} \ill{} \ill{} dans le \uncertain{noyau} ${}_n\Gamma$ est un sous-groupe \uncertain{ouvert et fermé}. Soit $\Gamma$ un $k$-groupe, loc.\ \uncertain{fini} sur $k$. [La \uncertain{suite} des ${}_n\Gamma$ est \struck{\ill{} et il existe un plus grand \ill{}} \add{\ill{} \ill{}}

\marginal{\uncertain{identité} cf plus bas}

\textbf{Proposition} \quad \struck{Il existe} \uncertain{La réunion} des \ill{} \ill{}, \ill{} ${}_n\Gamma$, est un \ill{} \ill{} \ill{} \ill{} $\Gamma^\tau$ de $\Gamma$. \ill{} \ill{} $a \in$ l'élément $a$ de $\Gamma$ \uncertain{soit} dans $\Gamma^\tau$, il faut et il suffit qu'il existe $n$ tel que $na = e$.

\marginal{\uncertain{après} \ill{} \ill{} des \ill{} \uncertain{finis} \ill{} $\Gamma$ \ill{}}

Soit $a \in \Gamma$, \struck{\ill{}} si $a \in {}_n\Gamma$ \ill{} \ill{} $na = 0$. \uncertain{Réciproquement}, \ill{} $na = 0$, comme il existe \ill{} $n$ tel que $n\Gamma_{\uncertain{a}} = 0$, [\uncertain{cf}] il s'ensuit que \struck{\ill{}} $\{a\} \xrightarrow{\text{ind}} \Gamma \xrightarrow{\ n\ } \Gamma$

\page{204}
est nul, i.e.\ \struck{\ill{}} \add{$\{a\}$} est dans le \uncertain{schéma noyau} \ill{} \ill{}, \ill{} $a$, donc $\{a\}$ est \ill{}. Cela \ill{} bien que $\Gamma$ est loc.\ \uncertain{constant}, que \ill{} $\Gamma^\tau$ qui est \ill{} \ill{}. Cela est \ill{} immédiat, et résulte \ill{} \add{qu.-cpt} \ill{}.

\struck{N.B.} \textbf{Corollaire} \quad Pour \ill{} $S$ \struck{\ill{}} sur $k$, \add{\ill{}}

$\Gamma^\tau(S)$ = torsion de $\Gamma(S)$. \struck{Et la \uncertain{plus}}

\struck{Corollaire $\Gamma^\tau$ est le plus petit sous-groupe \ill{} \ill{} tel que $\Gamma/\Gamma^\tau$ soit séparable \ill{} \ill{}}

Si $G$ est un \struck{\ill{}} $k$-groupe \ill{} de type fini, on \ill{} $G^\tau$ l'image inverse de \ill{} $\Gamma^\tau$ par $G \to G/G_0 = \Gamma$ (resp.\ $G^\tau_{\uncertain{\varepsilon}}$)

\textbf{Proposition} \quad $G^\tau$ \uncertain{est} le plus petit sous-\ill{} de $G$ tel que $G/G_0$ soit
\emph{discret} \emph{séparable} [resp.\ discret \struck{séparable} \add{sans} \struck{torsion}]. (C'est \ill{} \ill{} \ill{} \ill{} [\ill{} ; \ill{} : \ill{} \ill{}] \ill{} \ill{} \ill{} \ill{} \ill{})
$(G/G_0)^\tau = \{e\}$, \struck{\ill{}} $G^\tau$ est dit le sous-groupe de torsion de $G$.

\marginal{à \uncertain{préciser}}

\page{205}
$G$ $k$-groupe loc.\ de type fini sur $k$,

$G_0$ composante connexe,

$\Gamma = G/G_0$.

$\Gamma \xrightarrow{\ \pi_n\ } \Gamma$ \qquad ${}_n\Gamma = \pi_n^{-1}(e)$ \quad c'est un sous-schéma \emph{ouvert} et \emph{fermé} [car $\Gamma$ \uncertain{discret} \uncertain{séparable}] \ill{} \ill{} \ill{}

(i) La suite des ${}_n\Gamma$ est loc.\ stationnaire \ill{} \ill{}, \struck{\ill{}} (\uncertain{trivial} [car $\Gamma$ est discret] \ill{} \ill{} \ill{}) sa réunion est un sous-schéma \emph{ouvert} \struck{\emph{fermé}} $\Gamma^\tau$ de $\Gamma$ [N.B. c'est \uncertain{vrai} pour \uncertain{tout} $k$-groupe loc.\ fini sur $k$, pas nécessairement séparable, mais c'est moins évident, le cas radiciel \ill{} \ill{} \ill{} \ill{} \ill{} \uncertain{spécial}]

(ii) Soit $a \in \Gamma$, alors $a \in {}_n\Gamma^\tau$ ssi \struck{il existe \ill{}} $\pi_n(a) = e$.
$a \in \Gamma^\tau \Leftrightarrow \exists n,\ \pi_n(a) = e \Leftrightarrow$ il existe un \ill{} \struck{tel que \ill{}}
\struck{(iii) Soit} quelque \ill{} \ill{} [\ill{} \ill{} car $\Gamma$ sép.\ sur $k$]
$\Gamma'$ \uncertain{fini} sur $k$ tel que $a \in \Gamma'$ $\Longleftrightarrow$ l'un des $a^{\uncertain{n}}$ est \emph{fini}.

(iii) Si $\Gamma$ est \uncertain{commutatif}, alors $\Gamma^\tau$ est le plus grand sous-groupe [\ill{}] \emph{fini} \ill{} \ill{}. C'est \ill{} le plus petit \ill{} tel que $\Gamma/\Gamma^\tau$ soit « sans torsion ».

(iiii) Soit $G^\tau$ l'image inverse de $\Gamma^\tau$. \struck{\uncertain{Alors}} \add{C'est} un sous-schéma \ill{} de $G$, et
$x \in G^\tau \Longleftrightarrow \exists n,\ \pi_n(x) \in G_0 \Longleftrightarrow$ Il \uncertain{existe} \ill{} \ill{} $H$ de $G$ \ill{} \ill{}
$\Longleftrightarrow \exists$ \ill{} \ill{} \ill{} \emph{de type fini} \ill{} $G$ \ill{} \ill{} \ill{}

\section{Critère de type fini}

\page{207}
\marginal{$\mu(s) + 2\alpha(s)$ \ill{} \uncertain{continuité}}

\textbf{Proposition} \quad Soit $f : X \to Y$ un morphisme loc.\ de \uncertain{présentation finie}, \ill{} \ill{} \ill{} \uncertain{séparé} \ill{} \ill{} \ill{} loc.\ \uncertain{noethériens}. \uncertain{On} \uncertain{suppose} \ill{} (ii) $f$ \struck{\ill{}} \ill{} (iii) les fibres de $f$ \add{\uncertain{propres} \ill{} \ill{} \ill{}} \struck{de type fini} sont de type fini. Alors $f$ \ill{} \ill{} \ill{} \uncertain{supposer} $Y$ \uncertain{affine}, \ill{} \ill{}.

\marginal{[\ill{} \ill{} $\uncertain{X \times_Y}$ \ill{} \ill{} \ill{} $f : X \to$ \ill{} \ill{}]}

Il suffit de prouver que $f$ \ill{} \uncertain{quasi-compact}, \ill{} $\alpha$ \ill{} \ill{} \ill{} \ill{} \uncertain{Y} \ill{} \ill{} ceci : \struck{il existe \ill{} \ill{} \ill{}} \add{Il existe} \ill{} \ill{} \ill{} $Y$ \uncertain{tel que} $X|U$ \ill{}
\ill{} \ill{} il \uncertain{revient} au même de prouver \ill{} \ill{} \ill{} \ill{} \ill{} \ill{} \ill{}

\uncertain{qu.\ cpt.} [Il existe \ill{} \ill{} \uncertain{ouvert} \uncertain{quasi-compact} $V$ de $X$ \ill{} \ill{} \ill{} $X_\eta$. Il existe \ill{}
\struck{de type fini \ill{}} \add{\emph{propre}} \struck{\ill{}}
$X'$ \ill{} \ill{} \ill{} \ill{} $Y$, \ill{} \ill{} \ill{} \ill{}
\[
  X'_\eta \simeq X_\eta = V_\eta .
\]
Si on \uncertain{restreint} \ill{} \ill{} \ill{} $Y$, on peut supposer que le dernier iso provient
d'un isomorphisme $X' \simeq V$. Donc \struck{$V_\eta$ \ill{} \ill{}} \ill{} \ill{} \uncertain{propre} sur $Y$, donc \uncertain{fermé} dans $X$ [\ill{} $X$ \ill{}]. \struck{Alors $X_\eta$ \ill{} \ill{} \ill{} \ill{}} \struck{\ill{} \ill{} $X$}
\struck{\uncertain{f}} donc $X - V$ est \uncertain{ouvert}. \struck{$X - V$ \ill{}}
\struck{\ill{} \ill{}} Si on avait $X - V \neq \emptyset$, $X - V$ dominerait $Y$,
\struck{\ill{}} ce qui est incompatible avec $X_\eta = V_\eta$.
\ill{} $(X - V)_\eta = \emptyset$. Donc $X = V$, d'où la conclusion.

\page{208}
\textbf{Corollaire} \quad Soit $Z$ une partie de $X$. \ill{} \uncertain{supposons} (i) $Z$ loc.\ constructible dans $X$ (ii) $f : X \to S$ \uncertain{séparé} \uncertain{loc.} de type fini (iii) \struck{$Z$ \ill{}} \ill{} $x \in Z$, et \ill{} \uncertain{générisation} de $x$, $\exists$ \ill{} $x' \in Z$ \ill{} $x$, \ill{} \ill{} \ill{} \struck{\ill{} \ill{} \ill{} $x' \in X$} (iv) les fibres de $Z$ sur $S$ \ill{} \emph{propres}. Alors $Z \to S$ est
qu.-cpt.

On \uncertain{peut} supposer $S$ \uncertain{noethérien}, \ill{} \ill{} \ill{} \add{si $S \neq \emptyset$} qu'il existe un \ill{} $U \neq \emptyset$ \ill{} $Z|U$ soit qu.-cpt. On \uncertain{peut} supposer $S$ \ill{} \ill{} et \ill{} $Z'$ de type fini sur $S$ et \ill{} \ill{}
\struck{il existe un \ill{} $Z'$ propre \ill{}}
\ill{} $Z'_\eta \to Z_\eta$, \ill{} $X$ contenant $Z_\eta$. \uncertain{Restreignant} $S$,
\ill{} \ill{} \ill{} $Z' \to V$, [\ill{} \ill{} que $Z \cap V = Z'$]
\ill{} \uncertain{immersion}], \ill{} $Z'$ \ill{} \uncertain{propre} sur $S$.
Alors $Z'$ est \uncertain{fermé} dans $X$, et \ill{} \ill{}
\[
  Z = (Z \cap V = Z') \cup (Z - Z'),
\]
\add{\uncertain{constructibles}} \ill{} \ill{} \uncertain{disjoints}. \ill{} \ill{}
\ill{} \ill{} \ill{} $Z_\eta = \emptyset$,
\ill{} \ill{} $Z = \emptyset$ …

\note{le bas de la page est vide ; l'argument s'arrête sur « $Z = \emptyset$ … »}

\section{Partie constructible de $\underline{\mathrm{Pic}}$}

\page{209}
\textbf{Proposition} \quad $X$ projectif sur $S$ \add{noethérien}, $f_*(\mathcal{O}_X) = \mathcal{O}_Y$ universellement, fibres primaires, $P = \underline{\mathrm{Pic}}_{X/S}$, soit $P'$ la partie de $P$ réunion des $(P_s)'$, où $P'_s$ est l'ensemble des éléments de $P_s^\tau$ image inverse des sous-groupes de $P_s^\tau / P_s^0$ formés des éléments d'ordre premier à \struck{\ill{}} \add{car $k(s)$}.
\add{Supposons que dans $P^\tau$ l'adhérence de tt pt soit quasi-compacte} \add{(\uncertain{ou} \uncertain{que les} $(P_s)^0$ soient propres sur $k(s)$).}
Alors $P'_S$ est une partie \add{\emph{constructible}} \struck{\ill{} \ill{} de $P$} \emph{quasi-compacte} de $P$, \struck{\ill{}}

\textit{Démonstration.} Il suffit de prouver \add{(S \uncertain{loc.} noeth.} \ill{} $Y$) qu'il existe un ouvert non vide $U$ de $Y$, tel que la partie de $P'_S$ au-dessus soit qu.-cpt. On peut supposer de \uncertain{plus} $Y$ affine intègre, soit $\eta$ son pt générique. \struck{Maintenant} $P'_\eta$ est un \uncertain{sous}-groupe \add{\uncertain{ouvert}} de $P_\eta$, de type fini sur $k(\eta)$. Il provient d'un groupe \add{$H$} de type fini sur un \uncertain{voisinage} de $\eta$, et on peut supposer \add{(\uncertain{restreignant} $Y$)} que ce dernier est un ouvert dans $P^\tau$, et \uncertain{est} un sous-\emph{groupe} \uncertain{ouvert} de $P^\tau$ [\uncertain{prendre} un ouvert qu.-cpt $V$ de $P^\tau$ contenant $P'_\eta$].
\page{210}
\marginal{$H$ est \uncertain{propre} sur $S$.}
\add{Comme $\overline{V}$ est quasi-compact, on voit que \uncertain{à condition de} \uncertain{remplacer} $Y$ \uncertain{par} \uncertain{une partie}, $H'$ est fermé dans $P^\tau$. Mêmes \uncertain{conclusions} \ill{}} \struck{Il} suffit de prouver que \struck{\ill{} \ill{}} \add{\uncertain{sous ces}} \add{conditions} \struck{\uncertain{restreignant} $Y$}, on a $H = P'$. En effet, il \struck{\ill{}} \ill{} les $P_s^0$. Il suffit donc, compte tenu que $P_s^0$ est divisible par les entiers premiers à la car.\ de $k(s)$, de prouver que \struck{les éléments de $P_s$} pour tout entier $n$ premier à la car.\ de $k(s)$, ${}_n(P_s) \subset H$. Or ${}_nP_s = \underline{H}^1(X_s, {}_n\mu)$, d'autre part on sait (\uncertain{cf}) que sur l'ouvert $S_n$ où $n$ est inversible, $\underline{H}^1(X_n/S_n, {}_n\mu)$ est un \struck{\ill{} \ill{}} schéma \uncertain{étale} sur $S$, donc tout pt de sa fibre en $s$ est spécialisation d'un point de la fibre générique, i.e.\ d'un pt de ${}_n(P_\eta) \subset P'_\eta \subset H$. Comme $H$ est fermé, on a donc ${}_n(P_s) \subset H$, cqfd.

\textbf{Corollaire} \quad Supposons que (i) \add{les $(P_s)^0$ soient propres sur} \ill{} dans $P^\tau$ \struck{\ill{}} l'adhérence d'un pt soit quasi-compacte (ii) les $P_s^\tau / P_s^0$ \add{soient d'ordre premier à la car.\ $k(s)$}, \struck{\ill{} \ill{} \ill{} propres sur $k(s)$}. Alors $P^\tau$ est de type fini sur $S$.

\page{211}
\textbf{Corollaire} \quad Soit $P^0 = \bigcup_{s \in S} P_s^0$, supposons $P^0$ \emph{fermé}. Alors $P^\tau$ est fermé, \struck{et de type fini sur $S$} \add{si} \struck{\ill{}} pour tout $s \in S$, $P_s^\tau / P_s^0$ est d'ordre premier à la car.\ de $k(s)$, \struck{\uncertain{Prenons} \ill{} $P'$ est fermé}

\section{Un exemple}

\page{213}
\[
  S = \operatorname{Spec} k[t], \qquad X = \operatorname{Proj} k[t, x, y, z]/(xy - tz^2)
\]
\quad graduée en $x, y, z$ degré 1, $t$ degré 0.

\textbf{Lemme a)} \quad $(xy - tz^2)$ est un idéal premier dans $k[t, x, y, z]$, et l'anneau $B = k[t, x, y, z]/(xy - tz^2)$ est \uncertain{absolument} \uncertain{simple} sur $k$ [donc \uncertain{régulier}] \uncertain{sauf} \struck{\ill{}} en l'idéal \uncertain{d'irrégularité} \ill{}

\textbf{Cor} \quad $X$ est \struck{\ill{}} \uncertain{régulier}, et $f : X \to S$ est plat.

\textbf{Lemme b)} \quad Soit $a$ le pt $t = 0$ de $S$, $b$ \add{$\in U = S - \{a\}$} \add{\uncertain{autre} point.} \struck{le pt générique}. Alors (i) $f^{-1}(b)$ est isomorphe à $\mathbf{P}^1_{k(b)}$ [de \uncertain{façon} plus précise, $X|U \simeq \mathbf{P}^1_U$ \struck{\ill{} fibré} \struck{projectif \ill{} à \ill{} fibré vectoriel}]
(ii) $f^{-1}(a)$ est isomorphe \struck{\ill{}} à $\mathbf{P}^1_k \vee \mathbf{P}^1_k$

\textbf{Corollaire 1} \quad \add{\uncertain{suite} exacte} $\left[0 \to \mathcal{O}_{X_a} \to \mathcal{O}_{\widetilde{X}_a} \to k \to 0\right]$

\textbf{Corollaire 2} \quad $H^i(X_\xi, \mathcal{O}_{X_\xi}) = 0$ pour $i > 0$, $\xi \in S$, d'où \ill{}

\textbf{Corollaire 3} \quad \struck{Deux} \add{Soit} \struck{\uncertain{faisceaux} inversibles \ill{}} $S' = \operatorname{Spec}(A')$, $A'$ \ill{} local \uncertain{épointé} sur $k[t]$, $X' = X \times_S S'$. Alors
\[
  (*) \qquad H^1(X', \mathcal{O}^*_{X'}) \to H^1(X'_0, \mathcal{O}_{X'_0})
\]
est injectif ($X'_0$ étant la fibre \add{\uncertain{spéciale}} de $X'$) \struck{\uncertain{bijectif}} \add{\uncertain{bijectif}} \struck{et bijectif si $A'$ est complet.}

\note{dans $(*)$ on lit bien $\mathcal{O}_{X'_0}$ au but, non $\mathcal{O}^*_{X'_0}$}

\textbf{Corollaire 4} \quad Le deuxième membre de $(*)$, est \struck{\ill{}} isomorphe à $\mathbb{Z}$ si \struck{\ill{}} l'image du point fermé de $S'$ par $S' \to S$ est \struck{\ill{}} $\neq a$, et à $\mathbb{Z} \times \mathbb{Z}$ dans le cas contraire.

\marginal{\uncertain{Précisions}}

\page{214}
\textbf{Lemme c)} \quad Soient $Y_i$ $(i = 1, 2)$ les deux composantes de $f^{-1}(a)$, et $L_i$ les faisceaux inversibles associés \struck{\ill{}} aux diviseurs loc.\ principaux associés (cf \ill{}). Alors $L_i | Y_i$ est de degré $-1$, $L_i | Y_j$ de degré $1$ si $i \neq j$.

En \uncertain{effet}, on a $D_1 + D_2 = \struck{\ill{}}\operatorname{div}(t) \simeq 0$, d'où \struck{\ill{}} $L_{D_1 + D_2} | Y_1 \simeq 0$ i.e.
\[
  L_{D_1} | Y_1 \simeq - L_{D_2} | Y_1 .
\]
Or $L_{D_2} | Y_1 = 1$ par un calcul immédiat [qu'il faudrait \uncertain{savoir} faire] d'où $L_{D_1} | D_1 = 1$
\note{la dernière égalité porte $L_{D_1}|D_1 = 1$ ; avec ce qui précède on attendrait $-1$}

\textbf{Lemme e} \quad Il existe deux sections $s_1, s_2$ de $X/S$, telles que \struck{\ill{}} $D'_i$ \uncertain{dans} les diviseurs \uncertain{associés} \ill{} $L_i = L_{D'_i}$ \ill{} \add{\uncertain{deg} $L_i | Y_i = 1$, $\deg L_i | Y_j = 0$ si $i \neq j$}.

\textbf{Corollaire} \add{\uncertain{(préliminaire)}} \ill{} \ill{} \ill{}
\[
  0 \to \mathbb{Z} \xrightarrow{\ \alpha\ } H^1(X, \mathcal{O}^*_X) \xrightarrow{\ \beta\ } \mathbb{Z} \to 0
\]
où \struck{\ill{}} $\beta$ est \struck{\ill{}} \add{\uncertain{degré}} \ill{} la restriction à la fibre générique [ou \ill{} fibre spéciale \add{$f^{-1}(\xi)$, $\xi \neq a$} \ill{} …], et où $\operatorname{Ker} \beta$ est engendré par \struck{\ill{}} la classe $D_1$ [ou \uncertain{la} classe $D_2 = -$ classe $D_1$].

\textbf{Corollaire 2} \quad L'homomorphisme \add{\uncertain{de restriction}}
\[
  H^1(X, \mathcal{O}^*_X) \to H^1(X_a, \mathcal{O}^*_{X_a}) \simeq \mathbb{Z} \times \mathbb{Z}
\]
est \emph{bijectif}.
\note{une longue flèche, dans la marge gauche, relie ce dernier corollaire au « Lemme e »}

\page{215}
\struck{Démonstration du lemme d. Une section de $X$ sur $S$ \ill{} définie par un \ill{} \ill{} $k[t]$-algèbre \ill{} $\varphi : B \to A[z]$ ; \ill{} \ill{} \ill{} \ill{} $\varphi(z) = z$, $\varphi(x) = \xi$, \ill{} avec $\xi, \eta \in A$, $\xi\eta$ … soumis à $\xi\eta = t$}
\add{D'ailleurs on a} un ouvert \add{\uncertain{affine}} $V$ de $X$ isomorphe à $\operatorname{Spec} k[x, y, t]/(xy - t)$, et une section de $V$ sur $S$ est un $k[t]$-homomorphisme $k[x, y, t]/(xy - t) \to k[t]$, i.e.\ donnée par deux éléments $\xi, \eta$ de $k[t]$ tels que $\xi\eta - t = 0$.
1) $\xi = 1$, $\eta = t$ définit la section $s_1$. On regarde par la spécialisation $t = 0$ \ill{}
2) $\xi = t$, $\eta = 1$ …

\struck{Proposition} Démonstration du corollaire. \struck{Soit} \struck{Lemme e Sous les conditions, \uncertain{dans} \ill{} un diviseur sur $X$. \ill{} \ill{} à $X|U$ est \ill{} à un diviseur $D'_1$} [\ill{} \ill{}] \struck{\ill{} $H^1(U, \ldots)$}
\[
  H^1(X|U, \mathcal{O}^*_{X|U}) = H^1(\mathbf{P}^1_U, \mathcal{O}^*_{\mathbf{P}^1_U})
  = H^1(U, \mathcal{O}^*_U) \times \mathbb{Z} \simeq \mathbb{Z}
\]
compte tenu que $H^1(U, \mathcal{O}^*_U) = 0$ puisque $k[t][t^{-1}]$ est principal.
[N.B. \ill{} \ill{} \ill{} \ill{} \ill{} \ill{} \ill{} \ill{} \ill{} $H^1(S, P^1_S)$, \struck{\ill{}} qui ici se simplifie beaucoup grâce au fait que $S$ est \ill{}]

\page{216}
\struck{Donc \ill{} $\beta(\text{classe } D) = 0$, donc}
Donc le noyau de $\beta$ [\uncertain{ou} \uncertain{s'identifie} à $H^1(X, \mathcal{O}^*_X) \to H^1(U, \mathcal{O}^*_U)$] est le sous-groupe de $H^1(X, \mathcal{O}^*_X)$ engendré par les classes de $D_1$ et $D_2$, donc par celle de $D_1$, puisque $D_1 + D_2 \sim 0$. \struck{Donc} Or $\operatorname{cl}(D_1)$ \ill{} \ill{} un \ill{} \ill{} de $H^1(X, \mathcal{O}^*_X)$, grâce à lemme c). Le reste du corollaire 1 \add{et du lemme} \ill{} \ill{} \ill{} résulte \ill{} \uncertain{cor.\ 2}.

\textbf{Lemme e} \quad Sous les conditions du lemme b, corollaire 3, l'homom.\ $(*)$ est \emph{bijectif}.

\ill{} \ill{} \ill{} \ill{} \ill{} surjectif \ill{} \ill{} corollaire précédent.

On a \uncertain{ainsi} complété la \uncertain{structure} du Picard \uncertain{associé} : $f : X \to S$.

\textbf{Lemme f} \quad Le foncteur \add{$\mathcal{F}$} de Picard \struck{\ill{}} \struck{\ill{}} \uncertain{n'est pas} représentable par un préschéma sur $S$, loc.\ de type fini.
\add{Considérons en effet la fibre $f^{-1}(a)$, et l'élément de degré $(1, 1)$ [\uncertain{correspondant} : $L_2$] de $H^1(f^{-1}(a), \mathcal{O}^*_{f^{-1}(a)}) = \mathcal{F}(\operatorname{Spec} k(a))$.}
\note{la fin de la page est barrée d'un trait courbe ; les pages 216 (bas) et 217 (haut) sont cancellées par des diagonales}

\page{217}
\struck{Il lui correspondrait un point $\alpha$ de $\underline{\mathrm{Pic}}(X/S)$, rat./$k(a)$, et un anneau local $\mathcal{O}_\alpha$, tel que \ill{} \ill{} \ill{} \ill{} un homom.\ $\mathcal{O}_a \to \mathcal{O}_\alpha$, tel que pour \ill{} $\mathcal{O}_a$-algèbre \add{locale} \ill{} $A'$, \ill{} [$\mathcal{O}_a \to A'$ homom.\ local], on a $\mathcal{F}(A')_0 \simeq \operatorname{Hom}_{\mathcal{O}_a\text{-alg}}(\mathcal{O}_\alpha, A')$, où le \ill{} \ill{} $\mathcal{F}(A')$ désigne les \ill{} éléments de $\mathcal{F}(A')$ qui induisent l'élément de degré $(-1, 1)$ sur $\mathcal{F}(A'/\mathfrak{m}(A'))$. Comme il y en a exactement un, il faudrait donc que $\mathcal{O}_a = \mathcal{O}_\alpha$ \ill{}}

\struck{Considérons l'ensemble $\mathbb{Z} \times \mathbb{Z} \amalg \mathbb{Z} \times U$}
\uncertain{En} effet le préschéma \struck{\ill{}} $S^{\natural}_I$ \uncertain{suivant} [où $I$ est un ensemble d'indices arbitraire] : comme ensemble, il est identique à \struck{\ill{}} $U \times \{I\} \amalg I$ ; pour tout $i \in I$, soit $S_i = U \times \{I\} \cup \{i\}$, \uncertain{donc} on a
\[
  S_I = \bigcup_{i \in I} S_i , \qquad S_i \cap S_j = U \times \{I\}
\]
\ill{} \uncertain{correspondant} \uncertain{biunivoquement} avec $U$. On \uncertain{munit} chaque $S_i$ de la structure \add{de préschéma} \uncertain{définie} \ill{} \ill{} \ill{} \ill{} \ill{} de structure \ill{} \ill{} $S$, et \ill{} \uncertain{recolle}.
\note{$S_I$ est la droite $S$ dont l'origine est dédoublée en autant de points qu'il y a d'indices dans $I$}

\page{218}
Pour tout $n \in \mathbb{Z}$, soit $I_n \subset \mathbb{Z} \times \mathbb{Z}$ formé des couples $(p, q)$ \uncertain{tels que} $p + q = n$. Soit
\[
  \mathfrak{P} = \coprod_{n \in \mathbb{Z}} S_{I_n} .
\]
On a donc
\[
  \mathfrak{P} = \bigcup_{(p, q) \in \mathbb{Z} \times \mathbb{Z}} S_{(p,q)}
  \quad \text{où chaque } S_{(p,q)} \simeq S
\]
(réunion non disjointe). La fibre de $\mathfrak{P}$ en $a$ est $\simeq \mathbb{Z} \times \mathbb{Z}$, sa fibre en un $b \in U$ est $\simeq \mathbb{Z}$. On \uncertain{voit} que $\mathfrak{P}$ représente le foncteur Picard.
\note{l'argument s'arrête là ; le bas de la page est vide}

\section{Autre exemple}

\page{219}
\note{titre de sa main en haut à droite : « Autre exemple »}
$X$ projectif sur $S$ loc.\ noeth., $f_*(\mathcal{O}_X) = \mathcal{O}_S$.

Considérons $g : \hat{C}_X \to \hat{C}$ [\ill{} \uncertain{supposons} $X$ \ill{} \ill{} \ill{} \ill{} $\operatorname{Spec}(\underline{S})$, $\underline{S}$ engendré par $\underline{S}_1$.]

Alors $g_*(\mathcal{O}_{\hat{C}_X}) = \mathcal{O}_{\hat{C}}$ [\uncertain{vérif}]

d'où le fait que
\[
  H^1(\hat{C}, \mathcal{O}^*_{\hat{C}}) \to H^1(\hat{C}_X, \mathcal{O}^*_{\hat{C}_X})
\]
est \emph{injectif} [\uncertain{vérif} ; \ill{} \struck{\ill{}} loc.\ \uncertain{libres} quelconques]

Si \struck{\ill{}} $\xi$ est dans l'image, \ill{} [désignant par $\varepsilon : X \to \hat{C}_X$ la section nulle] \ill{} on a $\varepsilon^*(\xi) = 0$. Or [\uncertain{vérif}] $\xi = \varphi^*(\eta)(n)$ où
\struck{$\mathcal{O}_{\hat{C}_X}$ est le \ill{} \ill{}}
\ill{} est le \uncertain{faisceau} canonique sur $\hat{C}_X$ considéré comme fibré projectif sur $X$, d'où
\[
  \varepsilon^*(\xi) = (\varphi g)^*(\eta) \, \varepsilon^*(\mathcal{O}_{\hat{C}_X}(n)) = \eta ,
\]
\uncertain{donc} \ill{} $\eta = 0$, \uncertain{donc} $\xi = \mathcal{O}_{\hat{C}_X}(n)$, donc $\xi$ provient de $\mathcal{O}_{\hat{C}}(n)$. D'où
\[
  \boxed{H^1(\hat{C}, \mathcal{O}^*_{\hat{C}}) \simeq \mathbb{Z}, \quad \text{engendré par } \operatorname{cl}(\mathcal{O}_{\hat{C}}(1))}
\]

\page{220}
Soit $X$ \struck{\ill{}} projectif sur $S$ \uncertain{de} valuation discrète, \uncertain{définissons} une famille de \uncertain{quadriques} non singulières \uncertain{dégénérant} en un cône quadratique
\[
  X = \operatorname{Proj}\bigl(k[x, y, z, t]/(x^2 + y^2 + z^2 - \lambda t^2)\bigr)
  \quad \text{sur} \quad S = \operatorname{Spec}(\struck{k}\,\ill{})
\]
\struck{\ill{} \ill{}} \ill{} un corps [\ill{} \ill{} \ill{} \ill{} \ill{}, \uncertain{on} \uncertain{pourrait} \uncertain{prendre} $k = \mathbb{Z}$]. $X$ est \uncertain{normal} et \uncertain{irréductible} \struck{\ill{} sur $S$}. \struck{\ill{} fibres \ill{} irréductibles} Pour calculer $\underline{\mathrm{Pic}}_{X/S}$, il faut calculer $\operatorname{Pic}(X'/S')$ pour tout $S'$ \emph{local} sur $S$, i.e.\ $S' = \operatorname{Spec}(A)$, $A$ anneau local noethérien muni d'un $\lambda \in A$. \struck{\ill{}} 1) $\lambda \in \mathfrak{m}(A)$. Si $A$ \ill{} \ill{} \ill{}, on trouve $\mathbb{Z}$ (\add{engendré par $\mathcal{O}_X(1)$}, on démontre que cela reste vrai pour $A$ quelconque, grâce à $\boxed{H^1(X'_0, \mathcal{O}_{X'_0}) = 0}$. Si on \ill{} \ill{} \ill{} $\xi \notin \mathfrak{m}(A)$, \ill{} \ill{} \ill{} donne \struck{\ill{}} $\mathbb{Z}^2$, de générateurs $\alpha, \beta$. D'où \struck{\ill{}} \add{\uncertain{le} \uncertain{schéma} de Picard} \ill{} \ill{} \uncertain{partie} \ill{} de type fini sur $S$ \ill{} \ill{} $S$.

\marginal{\ill{} \ill{} \ill{} \uncertain{fibré} \ill{} \ill{} $\mathbf{P}^3_S$ \ill{} fibres \ill{}}
\marginal{\uncertain{à} \uncertain{mettre} \uncertain{ailleurs} : le schéma de Picard \uncertain{d'une} \ill{} est $\mathbb{Z}$, \ill{} \ill{} \uncertain{additif} \ill{} \ill{}}
\note{l'argument continue au-delà de ce lot}

\end{document}
